\documentclass[11pt]{article}
\usepackage[a4paper,margin=1.05in]{geometry}
\usepackage[T1]{fontenc}
\usepackage[utf8]{inputenc}
\usepackage{lmodern}
\usepackage{amsmath,amssymb,amsthm,mathtools}
\usepackage{enumitem}
\usepackage[hidelinks]{hyperref}
\usepackage{microtype}
\setlength{\emergencystretch}{3em}

\newtheorem{theorem}{Theorem}[section]
\newtheorem{proposition}[theorem]{Proposition}
\newtheorem{lemma}[theorem]{Lemma}
\newtheorem{corollary}[theorem]{Corollary}
\newtheorem{definition}[theorem]{Definition}
\newtheorem{remark}[theorem]{Remark}

\DeclareMathOperator{\Mod}{Mod}
\DeclareMathOperator{\I}{I}
\DeclareMathOperator{\J}{J}
\DeclareMathOperator{\K}{K}
\DeclareMathOperator{\Pack}{Pack}
\DeclareMathOperator{\diam}{diam}
\DeclareMathOperator{\rank}{rank}
\DeclareMathOperator{\ifr}{ifr}
\DeclareMathOperator{\ffr}{ffr}
\DeclareMathOperator{\Sp}{Sp}

\title{Invisible Flats and the Cost of Marked Normal Forms}
\author{Luca Blanchi}
\date{June 2026}

\begin{document}
\maketitle

\begin{abstract}
This note proves a coarse rate-distortion obstruction for reconstructing full
markings from finite-depth Torelli, Johnson, homological, level, or similar
partial data.  The general mechanism is simple and robust.  If an observable is
constant on a quasi-isometrically embedded abelian orbit
\(\mathbb Z^r\cdot x\), then every inverse normal form from the observable back
to a marked presentation has worst-case and average error \(\Omega(rM)\) on
boxes of side length \(M\).  The same invisible flat contains
\((M/R)^{\Omega(r)}\) many \(R\)-separated presentations, so reconstruction at
resolution \(R\) requires \(\Omega(r\log(M/R))\) auxiliary bits.

For mapping class groups this abstract obstruction is realized by standard
multitwist and bounded-support Johnson-filtration flats.  Unmarked curve data,
unmarked Jacobian data, and finite marking rigidifications leave invisible flats
of rank \(3g-3\) on the generic closed genus \(g\) locus.  Marked homological
data leave the Torelli kernel, which contains the separating multitwist flat of
rank \(2g-3\), and the same rank persists after adding the first Johnson
quotient.  For every fixed Johnson depth \(k\) on \(S_{g,1}\), the residual
kernel contains undistorted abelian subgroups of rank at least
\(c_k g-O_k(1)\).

Thus finite-depth Torelli or Johnson data may rigidify finite stabilizers or
determine an unmarked object, but they do not give low-cost marked normal forms.
The result is a presentation-theoretic lower bound: the large-scale geometry of
the fibre forces an arithmetic of missing marking information.
\end{abstract}

\section{Introduction}

A moduli problem often separates an object from a presentation of that object.
For a smooth curve \(C\) of genus \(g\), one may remember only the unmarked
curve, or one may remember a marking \(f:S_g\to C\).  Torelli-type theorems can
determine the unmarked object very strongly.  They do not automatically produce
a low-cost marked presentation.

This note quantifies the obstruction.  The relevant cost is not the difficulty
of recognizing the curve.  It is the coarse cost of choosing a full marking from
data that are constant on a large residual gauge group.

The guiding principle is:
\[
\text{large invisible flats in a fibre force large normal-form cost.}
\]
If an observable is constant on an undistorted orbit
\(\mathbb Z^r\cdot x\), then an inverse normal form must choose one output for a
box of mutually far marked presentations.  On an \(r\)-dimensional box of side
length \(M\), this forces error of order \(rM\), and it forces
\(r\log(M/R)\) bits of extra data to reconstruct the marking at resolution
\(R\).

The mapping class group provides natural high-rank invisible flats.  Pants
multitwists give rank \(3g-3\) flats in the full marking fibre.  Separating
multitwists give rank \(2g-3\) flats in the Torelli and Johnson-kernel fibres.
For deeper fixed Johnson levels, bounded-support local elements can be placed on
linearly many disjoint subsurfaces, giving invisible flats of rank \(\Omega_k(g)\).

\begin{theorem}[Informal main theorem]
Finite level, Torelli, first-Johnson, and fixed-depth Johnson observables do not
admit low-cost inverse normal forms to full markings.  More precisely, their
fibres contain undistorted abelian flats of rank linear in genus, and therefore
every inverse normal form has linear-genus coarse error and linear-genus missing
bit complexity on large boxes.
\end{theorem}

The theorem is compatible with classical Torelli.  Torelli determines the
unmarked curve from its principally polarized Jacobian on the smooth locus.  The
present result concerns the additional cost of recovering a marking.

\section{Coarse Marked Reconstruction}

Let \(X\) be a metric space of marked presentations, let \(\Gamma\) be a
finitely generated group acting on \(X\) by isometries or uniformly
quasi-isometric maps, and let
\[
q:X\to Z
\]
be an observable.  The fibre \(q^{-1}(q(x))\) is the set of presentations
indistinguishable from \(x\) by the observable.

\begin{definition}[Inverse normal form]
An inverse normal form for \(q\) is an arbitrary map
\[
N:q(X)\to X.
\]
No continuity, measurability, definability, or computability hypothesis is
imposed.  The lower bounds below therefore apply to all selectors.
\end{definition}

\begin{definition}[Invisible abelian quasi-flat]
Let \(A\cong \mathbb Z^r\) be generated by \(a_1,\ldots,a_r\).  The orbit
\(A\cdot x\) is a \(q\)-invisible abelian quasi-flat of rank \(r\) if
\[
q(ax)=q(x)
\qquad
\text{for all } a\in A
\]
and the orbit map
\[
\mathbb Z^r\to X,
\qquad
n=(n_1,\ldots,n_r)\mapsto a_1^{n_1}\cdots a_r^{n_r}x
\]
is a quasi-isometric embedding.  Equivalently, there are constants
\(\alpha>0\), \(\beta\ge0\) such that
\[
d_X(a^n x,a^m x)
\ge
\alpha |n-m|_1-\beta
\]
for all \(n,m\in \mathbb Z^r\).
\end{definition}

\begin{definition}[Invisible flat rank]
The invisible flat rank of \(q\) at \(x\) is
\[
\ifr_x(q)=
\sup\{r:\ q^{-1}(q(x))\text{ contains a }q\text{-invisible abelian quasi-flat of rank }r\}.
\]
The fibre flat rank \(\ffr_x(q)\) is defined similarly, but allows arbitrary
quasi-isometrically embedded copies of \(\mathbb Z^r\) in the fibre rather than
only abelian group orbits.  Always \(\ffr_x(q)\ge \ifr_x(q)\).
\end{definition}

\begin{proposition}[Coarse invariance of fibre flat rank]
Let \(q:X\to Z\) and \(q':X'\to Z'\) be observables.  Suppose that
\(F:X\to X'\) is a quasi-isometry and that
\[
q(x)=q(y)
\quad\Longrightarrow\quad
q'(F(x))=q'(F(y))
\]
on the region under consideration.  If \(q^{-1}(q(x))\) contains a
quasi-isometrically embedded copy of \(\mathbb Z^r\), then
\(q'^{-1}(q'(F(x)))\) also contains such a copy.
\end{proposition}

\begin{proof}
Compose the quasi-isometric embedding \(\mathbb Z^r\to q^{-1}(q(x))\) with
\(F\).  The result remains a quasi-isometric embedding, and the fibre-respecting
hypothesis puts its image inside one \(q'\)-fibre.
\end{proof}

\section{Invisible Flats and Rate-Distortion}

Let \(A\cong\mathbb Z^r\) be generated by \(a_1,\ldots,a_r\).  For
\[
n=(n_1,\ldots,n_r)
\]
write \(a^n=a_1^{n_1}\cdots a_r^{n_r}\), and put
\[
Q_M=[-M,M]^r\cap\mathbb Z^r.
\]

\begin{theorem}[Invisible flat rate-distortion]
Let \(A\cdot x\) be a \(q\)-invisible abelian quasi-flat, and suppose
\[
d_X(a^n x,a^m x)
\ge
\alpha |n-m|_1-\beta .
\]
Then every inverse normal form \(N:q(X)\to X\) satisfies
\[
\sup_{n\in Q_M}
d_X\bigl(N(q(x)),a^n x\bigr)
\ge
\alpha rM-\frac{\beta}{2}.
\]
Moreover, if \(n\) is uniformly distributed in \(Q_M\), then
\[
\mathbb E_{n\in Q_M}
d_X\bigl(N(q(x)),a^n x\bigr)
\ge
\alpha r\frac{M(M+1)}{2M+1}
-\frac{\beta}{2}.
\]
In particular, for \(M\ge1\), the average error is
\[
\Omega_{\alpha,\beta}(rM).
\]
\end{theorem}

\begin{proof}
Let \(n_+=(M,\ldots,M)\) and \(n_-=(-M,\ldots,-M)\).  Since \(A\cdot x\) is
invisible,
\[
q(a^{n_+}x)=q(a^{n_-}x)=q(x).
\]
The inverse normal form gives the same output for both points.  By the triangle
inequality, at least one of the two distances from \(N(q(x))\) is at least half
of \(d_X(a^{n_+}x,a^{n_-}x)\).  The quasi-isometric lower bound gives
\[
d_X(a^{n_+}x,a^{n_-}x)
\ge
2\alpha rM-\beta .
\]
This proves the worst-case estimate.

For the average estimate, pair \(n\) with \(-n\).  Again the two points have the
same observable value, so
\[
d_X(N(q(x)),a^n x)+d_X(N(q(x)),a^{-n}x)
\ge
d_X(a^n x,a^{-n}x)
\ge
2\alpha |n|_1-\beta .
\]
Averaging over \(Q_M\) gives
\[
\mathbb E d_X(N(q(x)),a^n x)
\ge
\alpha\,\mathbb E |n|_1-\frac{\beta}{2}.
\]
Since
\[
\mathbb E_{j\in[-M,M]\cap\mathbb Z}|j|
=
\frac{M(M+1)}{2M+1},
\]
we have
\[
\mathbb E_{n\in Q_M}|n|_1
=
r\frac{M(M+1)}{2M+1}.
\]
The claim follows.
\end{proof}

\begin{corollary}[Packing and missing bits]
Under the hypotheses of the theorem, the set
\[
F_M=\{a^n x:n\in Q_M\}
\]
contains an \(R\)-separated subset of cardinality at least
\[
C_\alpha
\left(
\frac{M}{R+\beta+1}
\right)^r
\]
for \(M\gg R+\beta\).  Consequently, any reconstruction procedure which is given
\(q(x)\) and an auxiliary message of \(b\) bits, and which must reconstruct every
point of \(F_M\) up to error \(R\), requires
\[
b\ge
r\log_2\!\left(\frac{M}{R+\beta+1}\right)-O_\alpha(r).
\]
\end{corollary}

\begin{proof}
Choose a coordinate subgrid of \(Q_M\) with spacing
\[
s=\left\lceil\frac{2R+\beta+1}{\alpha}\right\rceil .
\]
Distinct points in this subgrid have \(\ell^1\)-distance at least \(s\), hence
their orbit points are \(R\)-separated.  The subgrid has cardinality
\[
\gtrsim_\alpha
\left(
\frac{M}{R+\beta+1}
\right)^r .
\]
For the bit bound, a \(b\)-bit message gives at most \(2^b\) possible outputs over
the fixed observable value.  These outputs must form an \(R\)-net for \(F_M\),
so \(2^b\) is at least the packing number at a comparable radius.  Taking
logarithms proves the estimate.
\end{proof}

\begin{remark}[Flat-local sharpness]
On the invisible flat itself the bit estimate is sharp up to \(O(r)\): cover
the coordinate box \(Q_M\) by \(\ell^1\)-balls of radius comparable to \(R\) and
transmit the index of the ball.  This gives an upper bound
\[
r\log(M/R)+O(r)
\]
for reconstructing points on the flat.  The full fibre may have richer
nonabelian geometry, so the flat-local upper bound need not describe the entire
reconstruction problem.
\end{remark}

\section{Mapping Class Group Input}

Let \(S_g\) be a closed oriented surface of genus \(g\), and let \(S_{g,1}\) be
a genus \(g\) surface with one boundary component fixed pointwise.  Write
\[
\Gamma_g=\Mod(S_g),
\qquad
\Gamma_{g,1}=\Mod(S_{g,1}).
\]
Let \(\mathcal M(S)\) denote the marking graph.  The action of \(\Mod(S)\) on
\(\mathcal M(S)\) is geometric, so \(\mathcal M(S)\) is quasi-isometric to the
mapping class group with any word metric.

\begin{proposition}[Multitwist flats]
If \(\gamma_1,\ldots,\gamma_r\) are pairwise disjoint, pairwise non-isotopic
essential simple closed curves, then the twists
\[
T_{\gamma_1},\ldots,T_{\gamma_r}
\]
generate a subgroup \(\mathbb Z^r\), and its orbit in the marking graph is
quasi-isometrically embedded.
\end{proposition}

\begin{proof}
The twists commute and are independent because the curves are disjoint and
non-isotopic.  The Masur-Minsky distance formula detects the exponent of
\(T_{\gamma_i}\) in the annular projection to the annulus around \(\gamma_i\).
The annular contributions for disjoint curves add, giving a lower bound
\[
d_{\mathrm{mark}}\!\left(
x,T_{\gamma_1}^{n_1}\cdots T_{\gamma_r}^{n_r}x
\right)
\ge
c\sum_i |n_i|-C.
\]
The reverse inequality is immediate from applying the twists successively.
\end{proof}

We use the following standard rank inputs.  The maximal rank of a free abelian
subgroup of \(\Mod(S_g)\) is \(3g-3\), realized by a pants multitwist subgroup
\cite{BLM,FLM}.  The Torelli group
\[
\mathcal I_g=\ker\bigl(\Mod(S_g)\to\Sp_{2g}(\mathbb Z)\bigr)
\]
contains separating twists, and Vautaw's rank theorem gives maximal abelian rank
\(2g-3\) in \(\mathcal I_g\), realized by separating multitwists
\cite{Vau}.  The Johnson kernel \(\mathcal K_g\), the kernel of the first
Johnson homomorphism on \(\mathcal I_g\), contains the subgroup generated by
separating twists \cite{Joh,Put}.

For fixed Johnson depth we use a bounded-support locality input.  Let
\(\mathcal J_{g,1}(k)\) be the \(k\)-th Johnson filtration term, with the
convention \(\mathcal J_{g,1}(1)=\mathcal I_{g,1}\).
For every fixed \(k\), the Church-Putman bounded-support theorem gives a
constant \(G_k\), independent of \(g\), such that
\(\mathcal J_{g,1}(k)\) is generated by elements supported on homologically
standard subsurfaces of genus at most \(G_k\) \cite{CP}.  Standard local
nontriviality constructions, for instance the pseudo-Anosov elements in deep
Johnson filtration terms constructed in \cite{MP}, give an infinite-order
mapping class \(\eta_k\) supported on a surface of genus \(h_k\), with
\(\eta_k\in\mathcal J(k)\), whose extensions by identity lie in
\(\mathcal J_{g,1}(k)\).

\begin{lemma}[Disjoint local Johnson elements give flats]
Fix \(k\).  Suppose \(\eta_k\in\mathcal J(k)\) is an infinite-order element
supported on a homologically standard subsurface \(W_k\) of genus \(h_k\).  If
\(W_{k,1},\ldots,W_{k,r}\) are pairwise disjoint homologically standard copies
of \(W_k\) in \(S_{g,1}\), then, after replacing the corresponding extensions by
positive powers if necessary, they generate an undistorted subgroup
\[
\mathbb Z^r\le \mathcal J_{g,1}(k).
\]
\end{lemma}

\begin{proof}
The supported mapping classes commute because their supports are disjoint.
After passing to powers they are pure on their supports.  A relation among the
extensions restricts on each \(W_{k,i}\) to a power of \(\eta_k\), hence all
exponents vanish because \(\eta_k\) has infinite order.

For undistortedness, each infinite-order pure mapping class has either a
twisting component detected by an annular projection or a pseudo-Anosov
component detected by a nonannular subsurface projection.  These detecting
subsurfaces lie in the disjoint supports, so the corresponding terms add in the
Masur-Minsky distance formula.  This gives the required linear lower bound in
the sum of the exponents.
\end{proof}

\section{Finite and Torelli-Level Observables}

The first applications are closed-surface statements on the generic smooth
locus.  Finite automorphism groups are ignored in the usual coarse sense: they
do not change the existence or rank of undistorted flats.

\begin{theorem}[Unmarked curve and Jacobian data]
For a generic closed genus \(g\ge3\) curve, the observable which remembers only
the unmarked curve has invisible flat rank \(3g-3\).  The same is true for the
observable which remembers only the unmarked principally polarized Jacobian.
Therefore every inverse normal form from either observable to a full marking has
worst-case and average error \(\Omega(gM)\) on boxes of side length \(M\), and
requires \(\Omega(g\log(M/R))\) auxiliary bits at resolution \(R\).
\end{theorem}

\begin{proof}
Forgetting the marking leaves the full marking fibre, coarsely modeled by
\(\Mod(S_g)\).  A pants decomposition gives an invisible multitwist flat of rank
\(3g-3\), and no abelian subgroup of \(\Mod(S_g)\) has larger rank.  On the
generic smooth locus, Torelli identifies the fibre of the unmarked Jacobian
observable with the same unmarked curve fibre up to finite stabilizers.  The
rate-distortion bounds follow from the general theorem.
\end{proof}

\begin{theorem}[No finite rigidification]
Let \(q:\mathcal M(S_g)\to Z\) be an observable whose dependence on the marking
factors through a finite quotient of \(\Mod(S_g)\).  Equivalently, suppose there
is a finite-index subgroup \(K\le\Mod(S_g)\) such that
\[
q(kx)=q(x)
\qquad
\text{for all } k\in K.
\]
Then
\[
\ifr(q)=3g-3.
\]
In particular, finite level structures may remove finite stabilizers, but they
do not reduce the coarse cost of recovering a full marking.
\end{theorem}

\begin{proof}
Let \(P\cong\mathbb Z^{3g-3}\) be a pants multitwist subgroup.  Since \(K\) has
finite index in \(\Mod(S_g)\), the intersection \(K\cap P\) has finite index in
\(P\), hence has the same rank.  It is invisible for \(q\) and remains
undistorted.  The upper bound \(3g-3\) is the maximal abelian rank of
\(\Mod(S_g)\).
\end{proof}

\begin{corollary}[Homological level]
For every \(m\ge2\), the observable recording homological level \(m\) data leaves
an invisible flat of rank \(3g-3\).  Indeed, the level subgroup has finite index,
and powers \(T_\gamma^m\) of the twists in a pants decomposition act trivially on
\[
H_1(S_g;\mathbb Z/m\mathbb Z).
\]
\end{corollary}

\begin{theorem}[Marked homological data and the first Johnson quotient]
The observable whose residual gauge is the Torelli group has invisible flat rank
\[
2g-3.
\]
The same rank remains after adding the first Johnson quotient, so that the
residual gauge is the Johnson kernel \(\mathcal K_g\).  Hence these data still
force \(\Omega(gM)\) coarse marked reconstruction error and
\(\Omega(g\log(M/R))\) auxiliary bits on large boxes.
\end{theorem}

\begin{proof}
The Torelli kernel contains the separating multitwist subgroup associated to a
maximal family of disjoint separating curves.  This gives an undistorted
invisible flat of rank \(2g-3\).  Vautaw's theorem gives the matching upper
bound for abelian subgroups of the Torelli group.

The Johnson kernel contains all separating twists, while it is contained in the
Torelli group.  Therefore the same lower and upper bounds apply after adding the
first Johnson quotient.
\end{proof}

\section{Fixed-Depth Johnson Data}

We now work on \(S_{g,1}\), where the Johnson filtration is most naturally
defined.  A depth-\(k\) Johnson observable is an observable
\[
q_k:\mathcal M(S_{g,1})\to Z_k
\]
whose invisible gauge contains \(\mathcal J_{g,1}(k)\).

\begin{theorem}[Linear invisible flats in fixed Johnson depth]
For every fixed \(k\ge1\), there are constants \(c_k>0\) and \(C_k\ge0\) such
that, for all sufficiently large \(g\), the group \(\mathcal J_{g,1}(k)\)
contains an undistorted abelian subgroup
\[
A_{g,k}\cong\mathbb Z^{r_{g,k}},
\qquad
r_{g,k}\ge c_k g-C_k .
\]
\end{theorem}

\begin{proof}
Choose the local infinite-order element \(\eta_k\) supported on a
homologically standard subsurface \(W_k\) of genus \(h_k\) described above.
Embed pairwise disjoint copies of \(W_k\) into \(S_{g,1}\).  The number of such
copies is at least
\[
\left\lfloor\frac{g}{h_k}\right\rfloor-O_k(1).
\]
Extend \(\eta_k\) by identity on each copy.  Naturality of the Johnson
filtration puts all these extensions in \(\mathcal J_{g,1}(k)\).  The preceding
disjoint-support lemma shows that suitable positive powers generate an
undistorted free abelian subgroup of the same rank.  Taking
\[
c_k=\frac{1}{2h_k}
\]
and increasing \(C_k\) if needed gives the displayed linear lower bound.
\end{proof}

\begin{corollary}[No fixed-depth marked reconstruction]
Let \(q_k\) be a depth-\(k\) Johnson observable on \(\mathcal M(S_{g,1})\).
Every inverse normal form
\[
N:q_k(\mathcal M(S_{g,1}))\to\mathcal M(S_{g,1})
\]
has worst-case and average error
\[
\Omega_k(gM)
\]
on boxes of side length \(M\).  Its invisible fibre has packing entropy at least
\[
(M/R)^{\Omega_k(g)}
\]
in the coarse regime \(M\gg R\), and reconstruction at resolution \(R\) requires
\[
\Omega_k(g\log(M/R))
\]
auxiliary bits.
\end{corollary}

\begin{proof}
The invisible gauge of \(q_k\) contains \(\mathcal J_{g,1}(k)\), hence contains
the abelian quasi-flat from the theorem.  Apply the invisible flat
rate-distortion theorem and the packing corollary.
\end{proof}

\section{Presentation-Theoretic Interpretation}

The statements above are naturally presentation-theoretic.  The marked object is
a presentation, the mapping class group is a gauge group of presentation
changes, and the observable records partial data.  A normal form is a section of
the observable map.  Invisible flat rank measures how much coarse presentation
geometry remains inside a fibre.

This is the same principle as in no-free-degeneration results, but in a purely
coarse-geometric setting.  A normal form does not become low-cost merely because
the unmarked object is determined.  The missing marking information is measured
by large boxes inside the fibre.

The fixed-depth Johnson theorem is the main structural application.  It says
that making the observable deeper in a fixed Johnson tower does not remove the
linear-genus Euclidean sector of the marking fibre.  To make marked
reconstruction low-cost, the depth or the auxiliary marking information must
grow with the genus or with the desired scale.

\section{Boundary Conventions and Questions}

The sharp low-depth closed-surface statements use standard closed-surface rank
theorems.  The fixed-depth Johnson statement is written for \(S_{g,1}\) because
the Johnson filtration and the bounded-support theorem have their cleanest form
there.  Closed-surface variants follow whenever the observable's residual gauge
contains the images of the local subgroups constructed above under the relevant
capping or forgetting map.

Several refinements remain natural.
\begin{enumerate}[label=\textup{(\roman*)}]
\item Compute the optimal asymptotic constant
\[
\alpha_k=\limsup_{g\to\infty}\frac{1}{g}
\max\{r:\mathcal J_{g,1}(k)\text{ contains an undistorted }\mathbb Z^r\}.
\]
\item Develop a nonabelian rate-distortion theory for residual fibres, using
growth, hierarchical hyperbolicity, divergence, or asymptotic cones.
\item Compare information-theoretic lower bounds with explicit algorithms for
choosing markings from period, nilpotent, or level data.
\item Extend the invisible-flat framework to character varieties, local systems,
Higgs bundles, hyperkaehler moduli, and derived categories with large
autoequivalence groups.
\end{enumerate}

\begin{thebibliography}{99}

\bibitem[BBF]{BBF}
M. Bestvina, K. Bromberg, and K. Fujiwara,
\emph{Constructing group actions on quasi-trees and applications to mapping class groups},
Publications mathematiques de l'IHES 122 (2015), 1--64.

\bibitem[BM]{BM}
J. Behrstock and Y. Minsky,
\emph{Dimension and rank for mapping class groups},
Annals of Mathematics 167 (2008), 1055--1077.

\bibitem[BLM]{BLM}
J. S. Birman, A. Lubotzky, and J. McCarthy,
\emph{Abelian and solvable subgroups of the mapping class groups},
Duke Mathematical Journal 50 (1983), 1107--1120.

\bibitem[CP]{CP}
T. Church and A. Putman,
\emph{Generating the Johnson filtration},
Geometry \& Topology 19 (2015), 2217--2255.

\bibitem[FM]{FM}
B. Farb and D. Margalit,
\emph{A Primer on Mapping Class Groups},
Princeton University Press, 2012.

\bibitem[FLM]{FLM}
B. Farb, A. Lubotzky, and Y. Minsky,
\emph{Rank-one phenomena for mapping class groups},
Duke Mathematical Journal 106 (2001), 581--597.

\bibitem[Joh]{Joh}
D. Johnson,
\emph{An abelian quotient of the mapping class group \(\mathcal I_g\)},
Mathematische Annalen 249 (1980), 225--242.

\bibitem[MM]{MM}
H. Masur and Y. Minsky,
\emph{Geometry of the complex of curves II: Hierarchical structure},
Geometric and Functional Analysis 10 (2000), 902--974.

\bibitem[MP]{MP}
J. Malestein and A. Putman,
\emph{Pseudo-Anosov dilatations and the Johnson filtration},
Groups, Geometry, and Dynamics 10 (2016), 1409--1433.

\bibitem[Put]{Put}
A. Putman,
\emph{The Johnson homomorphism and its kernel},
Journal fuer die reine und angewandte Mathematik 735 (2018), 109--141.

\bibitem[Vau]{Vau}
W. R. Vautaw,
\emph{Abelian subgroups of the Torelli group},
Algebraic \& Geometric Topology 2 (2002), 157--170.

\end{thebibliography}

\end{document}
